\documentclass[conference]{IEEEtran}
\renewcommand{\ttdefault}{cmtt}
\usepackage{amsmath,amssymb}
\usepackage{graphicx}
\usepackage{booktabs}
\usepackage{cite}
\usepackage{url}
\usepackage{microtype}
\usepackage{array}

% A number that waits for a pass that has not finished; search for TBD before submission.
\newcommand{\TBD}[1]{\textbf{[#1]}}

\begin{document}

\title{Blind Single-Channel Separation of Coexisting 2G--5G Cellular Signals:\\
A Benchmark on 3GPP-Compliant Mixtures}

% If the review is double blind, replace this block by anonymous text.
\author{\IEEEauthorblockN{Hao Chen and Dayuan Tan}
\IEEEauthorblockA{Department of Electrical Engineering\\
\textit{contact details to be added}}}

\maketitle

% ============================================================
\begin{abstract}
% ============================================================
GSM, UMTS, LTE and 5G NR transmissions now share or neighbour the same spectrum, yet
there is little common ground for comparing methods that separate such signals from a
single received mixture. We study this problem on RFSS, a corpus of 100,000 mixtures of
two to four 3GPP-compliant sources, each passing through its own tapped-delay-line fading
channel and hardware impairments, and mixed co-channel or adjacent-channel. For the
two-source subset we give an exact definition of the per-source references, a
phase-sensitive permutation-invariant SI-SINR protocol with oracle references, and results
for FastICA, NMF and three deep separators (an STFT-domain BLSTM, DPRNN and Conv-TasNet),
each trained with three seeds and a learning rate chosen on validation data. On 7,526 held-out
test mixtures the deep models improve the SI-SINR over the input mixture by 5.6 to 6.1~dB,
which is 70\% to 76\% of the gain of an ideal-ratio-mask oracle (8.0~dB), while ICA and NMF
end below the input level. The ordering STFT-BLSTM, DPRNN, Conv-TasNet holds for every
matched seed pair, with mean differences of 0.2 to 0.5~dB. Absolute output quality remains low (about
+1.8~dB for the best model), so the task is far from solved. We state the limits of the study,
including that the data are simulated.
\end{abstract}

\begin{IEEEkeywords}
RF source separation, cellular coexistence, 3GPP, benchmark, permutation-invariant training.
\end{IEEEkeywords}

% ============================================================
\section{Introduction}
% ============================================================

Cellular generations coexist in shared and neighbouring bands~\cite{cabric2004implementation,andrews2014what}.
A receiver built for one standard sees the others as interference, and in dense deployments
that interference is both co-channel and adjacent-channel. Separating the individual waveforms
from one observation would let a receiver demodulate each standard in turn, and would give
spectrum monitoring and interference management a signal-level tool in addition to detection and
classification of the signals present~\cite{rajendran2018deep}.

Source separation has advanced fastest where shared benchmarks exist. Speech separation grew around
WSJ0-2mix~\cite{hershey2016deep}, permutation-invariant training~\cite{kolbaek2017multitalker}
and scale-invariant metrics~\cite{leroux2019sdr}, which gave Conv-TasNet~\cite{luo2019conv} and
DPRNN~\cite{luo2020dual} a common yardstick. For radio signals the nearest public benchmark, the
RF Challenge~\cite{lancho2024rfchallenge}, asks for the recovery of one known desired signal in the presence of one
interferer, so the receiver knows which signal to recover. The blind problem, with several unknown
standard-compliant sources, a varying source count and a permutation to resolve, has no such reference
point, and the widely used RadioML dataset contains single signals only~\cite{oshea2016radio}.

This paper takes a first step toward such a reference point for the two-source case. Our contributions are:
\begin{enumerate}
  \item A precise protocol on the RFSS corpus: how the per-source references are rebuilt from
    the stored waveforms, a phase-sensitive permutation-invariant SI-SINR, the gain over the input mixture as the main
    figure, and two oracle references that bound what is achievable.
  \item Results for FastICA, NMF and three deep separators trained under one recipe, with three seeds
    per deep model, matched-seed paired comparisons and intervals over 7,526 test mixtures.
  \item An account of what the numbers do and do not show: the gap to the oracle, the failure of the
    classical methods, the size of the model differences compared with seed variation, and the limits of a
    simulated, single-antenna benchmark.
\end{enumerate}

% ============================================================
\section{Related Work}
% ============================================================

\textit{RF datasets.} RadioML~\cite{oshea2016radio} supplies single-signal samples of synthetic modulations
and is a standard benchmark for modulation classification; it holds no mixtures and so no separation targets.
The RF Challenge~\cite{lancho2024rfchallenge} (ICASSP 2024) pairs one desired signal of known type with one
interferer and scores bit error rate and mean squared error. In our setting every source is unknown, the number
of sources is two to four, and outputs are matched to references by permutation.

\textit{Audio separation.} WSJ0-2mix~\cite{hershey2016deep}, WHAM!~\cite{wichern2019wham} and related corpora
established the protocol we borrow: per-source references, a permutation-invariant loss~\cite{kolbaek2017multitalker}
and SI-SNR~\cite{leroux2019sdr}. Conv-TasNet~\cite{luo2019conv} and DPRNN~\cite{luo2020dual} come from this
line of work; we test them on radio signals without domain-specific changes, next to an STFT-domain mask estimator.

\textit{Classical methods.} FastICA~\cite{hyvarinen2000ica} and NMF~\cite{fevotte2011nmf} are the usual blind baselines.
They use no labelled data, which is the reason to include them: they show how much a corpus-trained model adds.

% ============================================================
\section{The RFSS Corpus and Reference Signals}
\label{sec:data}
% ============================================================

\subsection{Signals, channels and impairments}

Each RFSS sample mixes two to four sources drawn from GSM (GMSK, $BT\!=\!0.3$)~\cite{3gpp45004}, UMTS
(W-CDMA, 3.84~Mcps)~\cite{3gpp25211}, LTE (OFDM, 1.4 to 20~MHz)~\cite{3gpp36211} and 5G NR
(OFDM, 30 or 120~kHz subcarrier spacing, 10 to 100~MHz)~\cite{3gpp38211}. Every source is generated from the
physical-layer specification of its standard at its native rate and then passes through its own SISO tapped-delay-line
channel (TDL-A to TDL-E of TR~38.901~\cite{3gpp38901}) with Jakes Doppler~\cite{jakes1994microwave} up to 700~Hz.
Each source then receives hardware impairments chosen at random, with parameter ranges informed by 3GPP base-station
conformance specifications~\cite{3gpp38104,3gpp36104,3gpp25104}: carrier frequency offset, I/Q imbalance, phase noise,
DC offset and a Rapp power-amplifier model~\cite{rapp1991effects}. About 20\% of sources are left clean, 30\% receive one
impairment and 50\% several. Sources are mixed either co-channel (all at baseband) or adjacent-channel (shifted by
standard-specific offsets), at per-sample SNRs between $-10$ and 40~dB, and white noise is added. A sample covers about 1~ms.
The corpus has 100,000 mixtures, split by index into 70\%, 15\% and 15\% for training, validation and test.
The two-source subsets used here hold 34,912, \TBD{val 2-src} and 7,526 mixtures. We use the single-antenna downlink
case throughout.

\subsection{References}

The released file stores, for each source, the waveform after the channel and the impairments, at its native rate
and before resampling, frequency shifting and power scaling. The reference for source $i$ is rebuilt from it with the
released function \texttt{build\_aligned\_references}: resample to the mixture rate, apply the adjacent-channel shift, apply the power
scaling. The noise is not stored; it equals the mixture minus the sum of the references. On 1,000 checked mixtures the
power of that residual agrees with the stored SNR to within 0.2~dB.

% ============================================================
\section{Task, Metric and References}
\label{sec:task}
% ============================================================

\subsection{Metric}

For a reference $\mathbf{s}$ and an estimate $\hat{\mathbf{s}}$ (both with mean removed) the SI-SINR is
\begin{equation}
\text{SI-SINR}(\hat{\mathbf{s}},\mathbf{s}) = 10\log_{10}\frac{\|\alpha\mathbf{s}\|^2}{\|\hat{\mathbf{s}}-\alpha\mathbf{s}\|^2},
\quad \alpha = \frac{\mathrm{Re}\{\hat{\mathbf{s}}^{H}\mathbf{s}\}}{\|\mathbf{s}\|^2}.
\label{eq:sisinr}
\end{equation}
The projection coefficient is a real scalar, so a constant phase rotation of the estimate reduces the score. This is
the phase-sensitive version appropriate for waveform recovery, and it is stricter than a magnitude or complex-gain
variant. Permutation invariance takes the output permutation that maximizes the mean over sources.
The \emph{gain} is the SI-SINR of the separated estimates minus that of the input mixture, where the input score is the mean over sources of
SI-SINR(mixture, source). The gain is our main figure because the input level varies with mixing mode and SNR (it averages
$-4.4$~dB over the test mixtures). All scores use 7,680-sample windows, drawn at a random offset that is fixed by a seed and the sample index.

\subsection{Oracle and classical references}

\emph{IRM oracle:} masks built from the reference magnitudes (2048-point STFT, hop 512) applied to the mixture STFT with the
mixture phase. It is a reference for magnitude masking, not a bound on every method.
\emph{Noise-limited oracle:} each reference plus the noise residual, the ceiling for perfect separation that leaves all additive noise in
every output.
\emph{FastICA:} on a Hankel embedding of the scalar mixture (window 256, hop 128), with the components matched to the references by the best permutation.
\emph{NMF:} Frobenius-norm NMF of the magnitude STFT with one component per source, Wiener-type masks and the mixture phase.

% ============================================================
\section{Deep Separators and Training}
\label{sec:methods}
% ============================================================

We train three models from scratch on the two-source training subset. The \emph{STFT-BLSTM} (7.36M parameters) predicts a complex ratio
mask per source from the log-magnitude STFT with a two-layer BiLSTM. \emph{DPRNN}~\cite{luo2020dual} (1.11M) uses
64 encoder filters, kernel 16 and six dual-path layers. \emph{Conv-TasNet}~\cite{luo2019conv} (2.52M) uses 256 filters,
kernel 16 and 24 temporal blocks. All use permutation-invariant training with a negative SI-SINR loss, Adam, gradient-norm clipping at 1.0,
batch size 8, cosine annealing to $10^{-5}$ over 10 epochs, and random 7,680-sample crops redrawn each epoch.
The peak learning rate is chosen per model on validation data: $3\!\times\!10^{-4}$ for STFT-BLSTM and Conv-TasNet and $10^{-3}$ for DPRNN,
the choice resting on the higher worst-seed validation gain. The checkpoint with the lowest validation loss is evaluated;
for the seed that was also used to choose the learning rate, the last epoch was kept. Each model is trained with three seeds.
The three models differ in size and cost per epoch, so a difference between them is a difference between these configurations under this
recipe and not between architecture families as such.

% ============================================================
\section{Results}
\label{sec:results}
% ============================================================

\begin{table}[t]
\centering
\caption{Two-source test results: gain over the input mixture (dB, PI-SI-SINR), mean over three seeds (range of the three seeds), and for the
adjacent-channel and co-channel mixtures with SNR above 20~dB ($n\!=\!1{,}105$ and $756$). All test mixtures: $n\!=\!7{,}526$.}
\label{tab:main}
\resizebox{\columnwidth}{!}{%
\begin{tabular}{lrrrr}
\toprule
Method & Params & All & Adjacent & Co-channel \\
\midrule
ICA & -- & -13.60 & -16.01 & -16.19 \\
NMF & -- & -1.24 & -2.61 & -2.38 \\
STFT-BLSTM & 7.36M & +6.14 (+6.12 to +6.15) & +6.94 & +6.91 \\
DPRNN & 1.11M & +5.85 (+5.76 to +5.90) & +6.58 & +6.50 \\
Conv-TasNet & 2.52M & +5.61 (+5.42 to +5.72) & +6.21 & +6.26 \\
\midrule
IRM oracle & -- & +8.04 & +10.86 & +10.08 \\
Noise-limited oracle & -- & +10.16 & +21.28 & +21.13 \\
\bottomrule
\end{tabular}

}
\end{table}

Table~\ref{tab:main} gives the main result. The three deep models improve the SI-SINR by 5.6 to 6.1~dB over the input, which is 70\% to
76\% of the IRM-oracle gain of 8.0~dB. ICA and NMF end below the input level (by 13.6 and 1.2~dB): when the mixture contains
unknown, overlapping, filtered and impaired waveforms, the independence and spectral-shape assumptions of these methods are only roughly met; we did not isolate
which failure dominates. The absolute output score of the best model is only about
$+1.8$~dB, because the input mixtures average $-4.4$~dB.

\textit{Order and seed variation.} The seed-to-seed range is 0.03~dB for STFT-BLSTM, 0.14 for DPRNN and 0.30 for Conv-TasNet. In every matched seed pair
(same seed number, independently initialised) STFT-BLSTM exceeds DPRNN, by 0.26, 0.25 and 0.37~dB, and DPRNN exceeds Conv-TasNet, by 0.17, 0.22 and 0.33~dB; the
sign is the same in the adjacent-channel and co-channel bins. The paired 95\% intervals over test mixtures are about $\pm0.03$~dB wide. We therefore report the
order STFT-BLSTM $>$ DPRNN $>$ Conv-TasNet for this recipe, with two cautions. The pairing by seed number is nominal, so the evidence is the consistent sign in three pairs together with
the small seed ranges and not a paired test across seeds. And the models differ in parameters (7.36M, 1.11M, 2.52M) and in cost, so the comparison
does not isolate the domain in which the mask is estimated. The gap between DPRNN and Conv-TasNet (0.2 to 0.3~dB) is small in practice.

\begin{figure}[t]
\centering
\includegraphics[width=\columnwidth]{figures/fig_snr.pdf}
\caption{Gain over the input mixture by per-sample SNR bin on the two-source test mixtures (mean over three seeds for the deep models).
Bins hold 1,083, 1,904, 2,678, 1,499 and 362 mixtures from low to high SNR.}
\label{fig:snr}
\end{figure}

\textit{Dependence on SNR.} Figure~\ref{fig:snr} shows the gain by SNR bin. All three deep models improve steadily with SNR, from about 5 dB in the
lowest bin to 6.7 to 7.7~dB in the highest, while the gap to the IRM oracle widens from nothing to 4 to 5~dB. In the lowest bin the best models match or exceed the
oracle, which is a reference built on the mixture phase and not a limit on every method. At high SNR the noise-limited oracle shows that
a large margin remains. The learned separators do not approach it.

\textit{Robustness of the evaluation.} The headline numbers use one fixed random window per mixture. A pass on the first 7,680 samples of each mixture and
\TBD{two further random-window passes} give gains within \TBD{x} dB of these for every model and leave the order unchanged.

% ============================================================
\section{Limitations}
\label{sec:limits}
% ============================================================

\textit{Simulation.} RFSS is simulated. Standard-compliant generation, 3GPP channels and impairments make it a controlled benchmark in the way WSJ0-2mix is for speech,
but it contains no over-the-air captures, and we do not claim that a model trained on it works on a real receiver. Its value is comparison under identical conditions, and a
sim-to-real study is future work.
\textit{Scope.} Single antenna, downlink waveforms, about 1~ms per mixture, two sources in this paper. Three- and four-source mixtures are in the corpus; we leave them to later work.
\textit{Absolute quality.} The best separated output scores about $+1.8$~dB, which is not enough to demodulate and we report no bit error rate or error vector magnitude. The gain
should be read as progress on a hard task and not as a deployable result.
\textit{Metric.} SI-SINR with a real projection coefficient is phase sensitive; magnitude-only or phase-aligned scores would be higher.
\textit{Training budget.} Ten epochs under-train the models: a validation probe of a twenty-epoch STFT-BLSTM run gained 0.20~dB more in the all-samples bin.
\textit{Statistics.} Three seeds per model, a nominal seed pairing, and models that differ in size.
\textit{Test-set use.} Before the final evaluation, one early checkpoint was scored on 600 two-source test crops and ICA and NMF were run on the full test split; neither informed any setting.
The final evaluation used checkpoints and a recipe frozen on validation data. We ran a first-window pass and the random-window pass, read the first before fixing the window protocol, and report the
random-window pass, which matches the validation crops, as the headline. Differences between the passes are given above.

% ============================================================
\section{Conclusion}
% ============================================================

On two-source RFSS mixtures, learned single-channel separators recover 70\% to 76\% of the IRM-oracle gain while ICA and NMF fall below the input, and the three deep models rank
consistently across seeds with mean margins of 0.2 to 0.5~dB. The remaining gap to the noise-limited oracle and the low absolute output level make the task an open problem.
The corpus, reference-building code and evaluation scripts are released with \TBD{RFSS v1.1 link}. Larger source counts, over-the-air validation and a downstream demodulation metric are the next steps.

\bibliographystyle{IEEEtran}
\bibliography{icc}

\end{document}
